\documentclass[11pt,a4paper]{article}

%% PACHETE MATEMATICE
\usepackage{amsmath,amssymb,amsfonts}
\usepackage{amsthm}
\usepackage{mathpazo}
\usepackage{eulervm}
\usepackage{enumerate}
\usepackage{color}
\usepackage{graphicx}
\usepackage[all]{xy}
\usepackage{xcolor}
\usepackage{framed}
\usepackage[unicode]{hyperref}
\setcounter{MaxMatrixCols}{20}
\usepackage{multicol}
%\usepackage[romanian]{babel}


%% ASPECT
\linespread{1.05}
\usepackage[T1]{fontenc}
\usepackage[utf8x]{inputenc}
\usepackage[margin=2cm]{geometry}
% \usepackage[color]{showkeys}
\usepackage{anyfontsize}
\usepackage{titlesec}
\titleformat*{\section}{\large\bfseries}

\usepackage{fancyhdr}
\pagestyle{fancy}
\rhead{\small \color{gray} Notițe de Adrian Manea}
\lhead{\small \color{gray} M1-ACS, 2017---2018, M. Olteanu}


\usepackage[autostyle = false, style = german]{csquotes}
\usepackage[romanian]{babel}
\newcommand\qq{\enquote}

%% COMENZILE MELE
\renewcommand*{\proofname}{Dem.:}
\newcommand\bb{\mathbb}
\newcommand\dr{\mathrm}
\newcommand\kal{\mathcal}
\newcommand\fk{\mathfrak}
\newcommand\op{\oplus}
\newcommand\ot{\otimes}
\newcommand\ds{\displaystyle}
\newcommand\id{\indent}
\newcommand\nid{\noindent}
\newcommand\seq{\subseteq}
\newcommand\sneq{\subsetneq}
\newcommand\speq{\supseteq}
\newcommand\spneq{\supsetneq}
\newcommand\rar{\rightarrow}
\newcommand\Rar{\Rightarrow}
\newcommand\xrar{\xrightarrow}
\newcommand\lar{\leftarrow}
\newcommand\Lar{\Leftarrow}
\newcommand\xlar{\xleftarrow}
\newcommand\lrar{\leftrightarrow}
\newcommand\Lrar{\Leftrightarrow}
\newcommand\ideal{\trianglelefteq}
\newcommand\ai{\text{ a.\^{i}. }}
\newcommand\sk{(\textit{Schi\c{t}\u{a}}) }
\newcommand\rhup{\rightharpoonup}
\newcommand\rhdn{\rightharpoondown}
\newcommand\lhup{\leftharpoonup}
\newcommand\lhdn{\leftharpoondown}
\newcommand\vid{\emptyset}
\newcommand\wh{\widehat}
\newcommand\ol{\overline}
\newcommand\NN{\mathbb{N}}
\newcommand\RR{\mathbb{R}}
\newcommand\ZZ{\mathbb{Z}}
\newcommand\QQ{\mathbb{Q}}
\newcommand\CC{\mathbb{C}}

%% STILURILE MELE
\newtheoremstyle{dotless}{}{}{\itshape}{}{\bfseries}{:}{ }{}

\theoremstyle{dotless}
\newtheorem{theorem}{Teorem\u{a}}[section]
\newtheorem{proposition}{Propozi\c{t}ie}[section]
\newtheorem{lemma}{Lem\u{a}}[section]
\newtheorem{corollary}{Corolar}[section]
\newtheorem{exercise}{Exerci\c{t}iu}

\newtheoremstyle{dotlessdr}{}{}{}{}{\bfseries}{:}{ }{}
\theoremstyle{dotlessdr}
\newtheorem{example}{Exemplu}[section]
\newtheorem{definition}{Defini\c{t}ie}[section]
\newtheorem{remark}{Observa\c{t}ie}[section]




\begin{document}

\begin{center}
  {\large\textbf{Seminar 6 \\ Funcții de mai multe variabile}}
\end{center}

\vspace{1cm}

\section{Limite și continuitate}
\label{sec:limite-cont}

Ne vom ocupa acum de studiul funcțiilor de mai multe variabile. Pentru simplitate,
rezultatele vor fi enunțate pentru funcții de două variabile, i.e.\
$f : A \seq \RR^2 \to \RR$, însă ele se extind identic și pentru $n$ variabile, pentru
orice $n \in \NN^\ast$.

Începem cu studiul limitelor iterate și continuității, pentru a ne adresa apoi
derivabilității și derivatelor parțiale.

\begin{definition}\label{def:lim-iter}
  Fie $f : A \to \RR$, cu $A \seq \RR^2$ și fie $(x_0, y_0)$ un punct de acumulare pentru
  mulțimea $A$.

  Atunci $\ell = \ds\lim_{(x, y) \to (x_0, y_0)} f(x, y)$ dacă și numai dacă, pentru orice
  $\varepsilon > 0$, există $\delta_\varepsilon > 0$, cu proprietatea că:
  \[
    \big( |x - x_0| < \delta_\varepsilon \land | y - y_0 | < \delta_\varepsilon \big) \Rightarrow %
    |f (x, y) - \ell | < \varepsilon.
  \]

  De asemenea, pentru o funcție de două variabile, se pot considera și \emph{limitele iterate},
  adică:
  \[
    \lim_{x \to x_0} \lim_{y \to y_0} f(x, y), \quad %
    \lim_{y \to y_0} \lim_{x \to x_0} f(x, y).
  \]
  Aceste limite, dacă există, nu sînt neapărat egale. Dacă există \emph{limita globală}, definită
  mai sus și dacă există fiecare din limitele din iterațiile de mai sus, atunci există și
  limitele iterate, iar acestea sînt egale.

  Dacă limitele iterate nu sînt egale, limita globală nu există.

  Se poate întîmpla, însă, ca limitele iterate să fie egale, dar limita globală să nu existe.
\end{definition}

Continuitatea se studiază separat:

\begin{definition}\label{def:cont-mm}
  Fie $f : A \to \RR$, cu $A \seq \RR^2$ o funcție de două variabile. Dacă funcția $f(x, y_0)$
  este continuă în orice punct $x_0$, cu $(x_0, y_0) \in A$, atunci $f$ este \emph{continuă parțial}
  în raport cu variabila $x$.

  Similar se definește continuitatea parțială în raport cu $y$.
\end{definition}

Continuitatea funcției implică continuitatea parțială în raport cu fiecare variabilă. Reciproca
este falsă.

O altă variantă specifică de continuitate este:
\begin{definition}\label{def:uc}
  Funcția $f$ se numește \emph{uniform continuă} pe $A$ dacă $\forall \varepsilon > 0$, există
  $\delta_\varepsilon > 0$ astfel încît $\forall (x_1, y_1), (x_2, y_2) \in A$:
  \[
    \big( | x_1 - x_2 | < \delta_\varepsilon \land | y_1 - y_2 | < \delta_\varepsilon \big) %
    \Rightarrow | f(x_1, y_1) - f(x_2, y_2) | < \varepsilon.
  \]
\end{definition}

Cu aceasta, avem următorul rezultat:

\begin{theorem}\label{thm:uc-c}
  Orice funcție uniform continuă pe $A$ este continuă pe $A$.

  O funcție continuă pe o mulțime închisă și mărginită (caz în care se numește \emph{compactă})
  este uniform continuă.
\end{theorem}

\begin{remark}\label{rk:siruri-limite}
  Limitele se mai pot calcula folosind legătura cu funcțiile. Amintim, $\ds\lim_{x \to x_0} f(x) = \ell$
  dacă și numai dacă, pentru orice șir $(a_n)$, cu limita $x_0$, avem $f(a_n) \to \ell$.

  Similar, putem folosi șiruri pentru doar una dintre variabilele unei funcții atunci cînd
  calculăm limite iterate.
\end{remark}

%%%%%%%%%%%%%%%%%%%%%%%%%%%%%%%%%%%%%%%%%%%%%%%%%%%%%%%%%%%%%%%%%%%%%%

\section{Exerciții: Limite și continuitate}
\label{sec:ex-lim-cont}

1. Să se arate, folosind definiția, că următoarele funcții nu au limită în origine:
\begin{enumerate}[(a)]
\item $f(x, y) = \frac{2xy}{x^2 + y^2}, (x, y) \neq (0, 0)$;
\item $f(x, y) = \frac{x^2 + 2x}{y^2 - 2x}, (x, y) \neq (0, 0)$.
\end{enumerate}

\emph{Soluție:} Vom folosi definiția cu șiruri, punînd în evidență șiruri de puncte,
convergente către $(0, 0)$.

(a) Folosim șiruri de puncte care se află pe drepte care trec prin origine. Aceste drepte
au ecuații de forma $y = mx$. Dacă funcția ar avea limită în origine, ar trebui ca:
\[
  \lim_{(x, y) \to (0, 0)} f(x, y) = \lim_{x \to 0, y = mx} \frac{2x(mx)}{x^2 + x^2 m^2} = %
  \frac{2m}{1 + m^2}.
\]
Cu această limită depinde de șirul folosit (de $m$, mai exact), rezultă că limita nu
există, în general.

(b) Calculăm, similar:
\[
  \lim_{(x, y) \to (0, 0)} f(x, y) = %
  \lim_{x \to 0, y = mx} \frac{m^2 x^2 + 2x}{m^2 x^2 - 2x} = -1.
\]

Însă, dacă luăm $y^2 = px$, adică $(x, y) \to (0, 0)$ pe parabola care trece prin origine,
cu $p \in \RR - \{2 \}$, obținem:
\[
  \lim_{x \to 0, y^2 = px} \frac{px + 2x}{px - 2x} = \frac{p + 2}{p - 2}.
\]
Rezultă că, pentru șiruri diferite, obținem limite diferite, deci funcția nu are limită
în origine.

\vspace{1cm}

2. Să se calculeze limitele:
\begin{enumerate}[(a)]
\item $\ds\lim_{(x, y) \to (0, 0)} \frac{xy}{\sqrt{xy + 1} - 1}$;
\item $\ds\lim_{(x, y) \to (0, 2)} \frac{\sin xy}{x}$;
\item $\ds\lim_{(x, y) \to (\infty, \infty)} \frac{x + y}{x^2 + y^2}$.
\end{enumerate}

\emph{Soluție:} (a) Calculăm direct:
\begin{align*}
  \lim_{(x, y) \to (0, 0)} \frac{xy}{\sqrt{xy + 1} - 1} &= \lim_{u \to 0} \frac{u}{\sqrt{u + 1} - 1} \\
                                                        &= \lim_{u \to 0} \frac{(\sqrt{u + 1} + 1)u}{u}\\
                                                        &= 2.
\end{align*}

(b) Similar:
\[
  \lim_{(x, y) \to (0, 2)} \frac{\sin xy}{x} = \lim_{(x, y) \to (0, 2)} \frac{\sin xy}{xy} \cdot y = 2.
\]

(c) Pentru $x, y > 0$, avem:
\[
  0 < \frac{x + y}{x^2 + y^2} = \frac{x}{x^2 + y^2} + \frac{y}{x^2 + y^2} \leq \frac{x}{x^2} + \frac{y}{y^2}.
\]

Cum $\ds\lim_{(x, y) \to (\infty, \infty)} \Big( \frac{1}{x} + \frac{1}{y} \Big) = 0$, rezultă
că și limita inițială este nulă.

\vspace{1cm}

3. Să se arate că funcția:
\[
  f(x, y) = \begin{cases}
    \frac{xy^2 + \sin(x^3 + y^5)}{x^2 + y^4}, & (x, y) \neq (0, 0) \\
    0, & (x, y) = (0, 0)
  \end{cases}
\]
este continuă parțial în origine, dar nu este continuă în acest punct.

\emph{Soluție:} Considerăm variabilele separat. Funcția:
\[
  f(x, 0) = \begin{cases}
    \frac{\sin x^3}{x^2}, & x \neq 0 \\
    0, & x = 0
  \end{cases}
\]
este continuă în origine, deoarece $\ds\lim_{x \to 0} f(x, 0) = 0 = f(0, 0)$.

Similar, funcția $f(0, y)$ este continuă în $y = 0$.

Dar funcția $f(x, y)$ nu este continuă în origine, deoarece avem:
\begin{align*}
  \lim_{(x, y) \to (0, 0)} f(x, y) &= \lim_{x \to 0, y^2 = px} \frac{px^2 + \sin(x^3 + p^\frac{5}{2} x^{\frac{5}{2}})}{x^2 + p^2 x^2} \\
                                   &= \lim_{x \to 0, y^2 = px} \Bigg[ \frac{p}{1 + p^2} + %
                                     \frac{\sin(x^3 + p^{\frac{5}{2}} x^{\frac{5}{2}})}{x^3 + p^{\frac{5}{2}} x^{\frac{5}{2}}} \cdot %
                                     \frac{x^3 + p^{\frac{5}{2}} x^{\frac{5}{2}}}{x^2 + p^2 x^2} \Bigg] \\
                                   &= \frac{p}{1 + p^2},
\end{align*}
care depinde de $p$.

\vspace{1cm}

4. Să se arate că funcția:
\[
  f(x, y) = \begin{cases}
    \frac{2xy}{x^2 + y^2}, & (x, y) \neq (0, 0) \\
    0, & (x, y) = (0, 0)
  \end{cases}
\]
este continuă în raport cu fiecare variabilă în parte, dar nu este continuă în raport
cu ansamblul variabilelor.

\emph{Indicație:} Considerăm $y = b$, fixat. Avem $f(x, b)$ continuă. Similar, $f(a, y)$.

Pentru continuitatea globală în origine, luăm $y = mx$ și obținem o limită care depinde de $m$.

\vspace{1cm}

5. Să se arate că funcția:
\[
  f(x, y) = \begin{cases}
    \frac{xy + x^2 y \ln |x + y|}{x^2 + y^2}, & (x, y) \neq (0, 0) \\
    0, & (x, y) = (0, 0)
  \end{cases}
\]
este continuă parțial în origine, dar nu este continuă în raport cu ansamblul variabilelor în acest punct.

\emph{Indicație:} Limitele iterate în origine sînt ambele nule, indiferent de ordine.

Însă pentru $x \to 0$ și $y = mx$, limita globală depinde de $m$.

%%%%%%%%%%%%%%%%%%%%%%%%%%%%%%%%%%%%%%%%%%%%%%%%%%%%%%%%%%%%%%%%%%%%%%

\section{Derivate parțiale. Diferențiabilitate}
\label{sec:der-part}

Studiem acum proprietatea de derivabilitate a unei funcții de mai multe variabile. Deoarece
o funcție de două variabile, de exemplu, reprezintă o suprafață în spațiul tridimensional,
trebuie să ne orientăm după anumiți versori pentru a putea face operații de analiză.

De aceea, avem:

\begin{definition}\label{def:der-vers}
  Spunem că funcția $f$ este \emph{derivabilă în punctul $a$ după versorul $s$} dacă funcția
  reală $g : (-r, r) \to \RR, g(t) = f(a + ts)$ este derivabilă în punctul $t = 0$.

  În acest caz, numărul real $\frac{df}{ds}(a) = g'(0)$ se numește \emph{derivata lui $f$ %
    după versorul $s$ în punctul $a$} sau \emph{derivata lui $f$ după direcția $s$}.

  Echivalent, ca în cazul 1-dimensional, putem defini:
  \[
    \frac{df}{ds}(a) = \lim_{t \to 0} \frac{f(a + ts) - f(a)}{t}.
  \]
\end{definition}

Derivata parțială se definește prin limita după versorul $e_k$, din baza canonică
a spațiului $\RR^n$, unde funcția este considerată $f : \RR^n \to \RR$.

Dacă există derivata după $e_k$, ea se notează $\frac{\partial f}{\partial x_k}$
și se numește \emph{derivabilă parțial în raport cu variabila $x_k$}.

Dacă este derivabilă în raport cu toate derivatele, atunci spunem că $f$ este derivabilă
parțial pe domeniul de definiție.

Funcția $f$ se numește \emph{de clasă $C^1$ pe domeniul de definiție} dacă toate
derivatele parțiale ale sale există și sînt continue pe $D$.

Următoarea noțiune va fi de folos pentru a studia monotonia unei funcții de mai
multe variabile:
\begin{definition}\label{def:grad}
  Presupunem că $f : D \seq \RR^n \to \RR$ și că $f \in C^1(D)$. Atunci, pentru
  orice $a \in D$, vectorul din $\RR^n$ definit prin:
  \[
    \mathrm{grad}_a(f) = (\nabla f)(a) = \Big( \frac{\partial f}{\partial x_i}(a) \Big)_i
  \]
  se numește \emph{gradientul lui $f$ în punctul $a$}.

  De asemenea se mai definește numărul:
  \[
    \mathrm{div}_a(f) = (\nabla \cdot f)(a) = \sum_{i = 1}^n \frac{\partial f}{\partial x_i} (a),
  \]
  care se numește \emph{divergența} lui $f$.
\end{definition}

Fie $F : D \to \RR^m, F = (f_1, \dots, f_m)$, cu fiecare $f_i : \RR^n \to \RR$. Atunci:
\[
  J_F (a) = \Big( \frac{\partial f_i}{\partial x_j}(a) \Big)_{i, j}.
\]
Matricea se numește \emph{matricea jacobiană} a lui $f$, iar determinantul său
se numește \emph{jacobianul} aplicației.

Jacobianul se mai notează:
\[
  \det(J_f(a)) = \frac{D(f_1, \dots, f_m)}{D(x_1, \dots, x_m)}.
\]

Funcția $F$ se numește \emph{de clasă $C^1$ pe $D$} dacă fiecare componentă
a sa este de clasă $C^1$ pe domeniul de definiție.

Dată o funcție de mai multe variabile, se poate defini \emph{diferențiala} ei,
care joacă rol de \qq{derivată totală}, definită prin formula:
\[
  Df(a)(x_1, \dots, x_n) = \sum_{i = 1}^n \frac{\partial f}{\partial x_i}(a) dx_i.
\]

Ca în cazul funcțiilor de o singură variabilă, avem \emph{formula de derivare a %
  funcțiilor compuse} (numită și \emph{regula lanțului}, eng.\ ``chain rule''):
\[
  \frac{\partial f}{\partial x} = \frac{\partial f}{\partial y} \cdot \frac{\partial y}{\partial x}.
\]

Pentru diferențiale, se scrie în forma cunoscută:
\[
  D(G \circ F)(a) = DG(b) \circ DF(a),
\]
unde $A \seq \RR^n, B \seq \RR^m$, iar $A \xrar{F} B \xrar{G} \RR^k$, cu condiția
ca $F$ să fie diferențiabilă în $a \in A$ iar $G$ să fie diferențiabilă în $b = F(a)$.

Relația între matricele jacobiene este:
\[
  J_{G \circ F}(a) = J_G(b) \cdot J_F(a).
\]

Mai mult, aplicarea succesivă a operatorului de derivare parțială sau a celui de
diferențiere rezultă în derivate și diferențiale de ordin superior.

De exemplu:
\[
  \frac{\partial}{\partial x} \Big( \frac{\partial f}{\partial y} \Big) = %
  \frac{\partial^2 f}{\partial x \partial y}.
\]

Un rezultat util este:

\begin{theorem}[Teorema de simetrie a lui Schwartz]\label{thm:schwartz}
  Dacă $f \in C^2(D)$, atunci:
  \[
    \frac{\partial^2 f}{\partial x_j \partial x_k} (a) = %
    \frac{\partial^2 f}{\partial x_k \partial x_j} (a), \forall j, k, \forall a \in D.
  \]
\end{theorem}

Pentru diferențiala a doua, se definește o nouă matrice, numită \emph{matricea hessiană}:
\[
  H_f(a) = \Big( \frac{\partial^2 f}{\partial x_i \partial x_j} \Big)_{i, j}.
\]
Iar diferențiala a doua se scrie, atunci:
\[
  D^2f(a)(x_1, \dots, x_n) = \sum_i \sum_j \frac{\partial^2 f}{\partial x_i \partial x_j} (a) dx_i dx_j.
\]

Se poate defini un nou operator, pentru derivatele de ordinul al doilea, numit \emph{laplacian},
sau \emph{operatorul lui Laplace}:
\[
  \Delta f = \sum_i \frac{\partial^2 f}{\partial x_i^2}.
\]

%%%%%%%%%%%%%%%%%%%%%%%%%%%%%%%%%%%%%%%%%%%%%%%%%%%%%%%%%%%%%%%%%%%%%%

\section{Exerciții: Derivate parțiale}
\label{sec:ex-dp}

1. Fie $f : \RR^2 \to \RR, f(x, y) = 2x^3y - e^{x^2}$.

Folosind definiția, să se calculeze derivatele parțiale de ordinul întîi ale lui $f$
în punctele $(0, 0)$ și $(-1, 1)$.

\vspace{1cm}

2. Fie $f : \RR^3 \to \RR^2, f(x, y, z) = (x^2 - yz, y^2 - z^2)$.

Să se calculeze, folosind definiția, derivata după direcția $h = \big( \frac{1}{3}, \frac{2}{3}, \frac{2}{3} \big)$
a funcție $f$ în punctul $(x, y, z) = (1, 1, 2)$.

\vspace{1cm}

3. Fie $f : \RR^3 \to \RR^2, f(x, y, z) = (x^3 - y^3, x^3 + y^3 + z^3)$.

Să se calculeze derivata după direcția $h = \frac{1}{\sqrt{6}} (1, 1, -2)$ a lui $f$ în
punctul $(1, -1, 1)$.

\vspace{1cm}

4. Determinați gradientul, divergența  și laplacianul în punctul
$a = (1, 2)$, precum și matricea hessiană pentru funcțiile:
\begin{enumerate}[(a)]
\item $f(x, y) = 3\sin x + e^{\cos(x + y)}$;
\item $f(x, y) = x^3 y + 2x^2\cos y$;
\item $f(x, y) = \ln(x^2 + y^2) + 3xy$;
\item $f(x, y) = 2\sqrt{xy} + x^2y^3 - 3x + 2y$;
\item $f(x, y) = 5xe^y + x^2 + 5y^3$;
\item $f(x, y) = \frac{1}{\sqrt{x^2 + y^2}}$;
\item $f(x, y) = \sin(x^2y^3) - \arctan (x^3 y^2)$.
\end{enumerate}


\end{document}